\chapter*{Introduction}
\addcontentsline{toc}{chapter}{Introduction}

The bar complex of a chiral algebra carries a differential that
lives in~$\C$ and a coproduct that lives in~$\R$.

The differential $d_{\barB}$ extracts OPE residues along collision
divisors in the Fulton--MacPherson compactification $\FM_k(\C)$:
it records how operators interact when they collide in the
holomorphic variable~$z$.  The coproduct~$\Delta$ splits an
ordered sequence of tensor factors into two sub-sequences: it
records how an interval is cut at a point in the topological
variable~$t$.  A bar element of degree~$k$ is parametrized by
the product
\[
\FM_k(\C) \times \Conf_k(\R).
\]
The differential sees the first factor.  The coproduct sees the
second.  This product is the operation space of a two-colored
operad---the holomorphic--topological Swiss-cheese operad
$\SCchtop$---and the bar complex, equipped with both structures,
is an algebra over it.  That identification is the central theorem
of this volume.  To see it at work, consider the simplest chiral
algebra with nontrivial content.


\section*{The Heisenberg atom}

The Heisenberg algebra $\cH_k$ has a single generator~$J$, a
single OPE
\[
J(z)\,J(w) \;\sim\; \frac{k}{(z-w)^2},
\]
and no composite fields.  One generator, one pole order, no
composites: every step of the three-dimensional machine is
checkable by hand.

The bar complex is $\barB(\cH_k) = T^c(s^{-1}\C \cdot J)$, the
cofree tensor coalgebra on a single desuspended generator.  Its
differential is the coderivation whose cogenerator projection in
arity~$2$ is
\[
m_2(J, J) \;=\; k \cdot \mathbf{1},
\]
the OPE residue, with all higher operations vanishing:
$m_{k \geq 3} = 0$, because the Heisenberg OPE has no higher
singularities and there are no composite fields to produce them.
The coproduct is deconcatenation.  Together---differential from
$\FM_2(\C)$, coproduct from $\Conf_k(\R)$, no higher
interactions---the bar complex is a Swiss-cheese algebra, and a
formal one: the $\R$-factorization is exact.

On cohomology $H^\bullet(\cH_k, Q) = \C[J]$, the binary
operation~$m_2$ decomposes.  Its regular part is a commutative
product; its singular part is a $\lambda$-bracket:
\[
\{J \,{}_\lambda\, J\} \;=\; k\lambda.
\]
Together they form a $(-1)$-shifted Poisson vertex algebra.  The
Jacobi identity is vacuous (single generator); the Leibniz rule
is the derivation property of~$\partial$; skew-symmetry is the
antisymmetry of the OPE.  The entire physical content of this PVA
is in the $\lambda$-linear term---the level~$k$.

The Laplace transform of the $\lambda$-bracket produces the
classical $r$-matrix $r(z) = k/z^2$, and the quantum $R$-matrix
$R(z) = \exp(k/z)$---the exponential of a single pole---trivially
satisfies Yang--Baxter.

At genus~$1$, the Arnold relation acquires a defect from the
period matrix of the elliptic curve.  The bar differential
acquires curvature:
\[
d^2 \;=\; k \cdot \omega_1.
\]
The modular characteristic is $\kappa(\cH_k) = k$: the same
scalar that governs the OPE, the $\lambda$-bracket, the
$R$-matrix, and now the genus-$1$ obstruction.  The free energy
$F_1 = -(k/24)\,\log\,\eta(\tau)$ recovers the Dedekind
$\eta$-function.  Every datum of the Heisenberg
atom---bar differential, PVA, $R$-matrix, genus-$1$ curvature,
free energy---is controlled by this single parameter.  The
Rosetta Stone chapter (Section~\ref{sec:rosetta-stone}) verifies
that each identification instantiates the corresponding theorem of
Volume~I.


\section*{Presentations}

The Heisenberg computation exhibits a general phenomenon.  The
product $\FM_k(\C) \times \Conf_k(\R)$ that parametrizes bar
elements is the operation space of the two-colored Swiss-cheese
operad $\SCchtop$ (Definition~\ref{def:SC}), whose closed
(holomorphic) color has operations on $\FM_k(\C)$ and whose open
(topological) color has operations on
$\FM_k(\C) \times \Eone(m)$.  The bar complex is an algebra over
this operad (Theorem~\ref{thm:bar-swiss-cheese}).  The Recognition
Theorem (Theorem~\ref{conj:recognition}) characterizes which
$\Ainf$ chiral algebras carry this structure: any algebra
satisfying the four standing analytic hypotheses
\textup{(H1)--(H4)} is $\SCchtop$-structured.  The directionality
axiom---open inputs cannot produce closed outputs---is the
mathematical expression of bulk-to-boundary restriction.

Homotopy-Koszulity of~$\SCchtop$
(Theorem~\ref{thm:homotopy-Koszul}, proved via Kontsevich
formality and transfer from the classical Swiss-cheese operad)
ensures that the bar-cobar adjunction is a Quillen equivalence.
The bar complex therefore \emph{presents} the Swiss-cheese
algebra: it determines the algebra up to quasi-isomorphism, as a
free resolution determines a module.

The word \emph{presents} is precise and has a classical model.
The Steinberg variety
$\widetilde{\mathcal{N}} \times_{\mathcal{N}}
\widetilde{\mathcal{N}}$---the fiber product of the Springer
resolution with itself over the nilpotent cone---carries a
convolution product from composition of correspondences and a
stratification from the Bruhat decomposition.  The convolution
algebra is the Hecke algebra~$H(W,q)$: not defined by generators
and relations, but recovered from the geometry of the
correspondence space.  The Kazhdan--Lusztig polynomials, the BGG
correspondence, the character formulas for Category~$\cO$---all
extracted from the perverse filtration on the Steinberg variety.
The bar complex is the chiral-algebraic Steinberg variety.

\begin{center}
\renewcommand{\arraystretch}{1.35}
\begin{tabular}{p{0.42\textwidth}p{0.42\textwidth}}
\textbf{Steinberg $\to$ Hecke} &
  \textbf{Bar $\to$ Swiss-cheese} \\
\hline
$\widetilde{\mathcal{N}} \times_\mathcal{N}
  \widetilde{\mathcal{N}}$ &
  $T^c(s^{-1}\bar{\cA})$ on
  $\FM_k(\C) \times \Conf_k(\R)$ \\[3pt]
Convolution &
  Bar differential $d_{\barB}$ \\[3pt]
Bruhat stratification &
  Deconcatenation $\Delta$ \\[3pt]
$H(W,q)$ &
  $\SCchtop$-algebra \\[3pt]
KL polynomials &
  Spectral $R$-matrix $R(z)$ \\[3pt]
Category~$\cO$ &
  Line operators $\simeq \cA^!\text{-mod}$ \\[3pt]
BGG reciprocity &
  Koszul duality: boundary $= \cA^!$
\end{tabular}
\end{center}

\noindent
In both columns, the presentation-by-correspondence is the
content of a Koszul duality theorem.  At genus~$g \geq 1$, the
parallel deepens: the bar complex acquires curvature
$d^2 = \kappa(\cA) \cdot \omega_g$ from the Hodge bundle, and the
Swiss-cheese algebra curves---just as the Hecke algebra deforms
with its parameter~$q$.

Each of Volume~I's five theorems acquires a three-dimensional
reading through this lens:

\begin{center}
\renewcommand{\arraystretch}{1.3}
\begin{tabular}{p{0.06\textwidth}p{0.25\textwidth}p{0.56\textwidth}}
\textbf{Thm} & \textbf{Volume~I} &
  \textbf{Three-dimensional reading} \\
\hline
\textbf{(A)} & Bar-cobar adjunction
  & Swiss-cheese presentation; Quillen equivalence
    (\S\ref{sec:foundations}, \S\ref{sec:bar_cobar}) \\[3pt]
\textbf{(B)} & Bar-cobar inversion
  & Raviolo VA descent; boundary VA as cobar
    (Appendix~\ref{sec:raviolo}) \\[3pt]
\textbf{(C)} & Complementarity
  & The Koszul triangle inherits the
    $(-1)$-shifted symplectic structure
    (\S\ref{sec:ht-bulk-boundary-line}) \\[3pt]
\textbf{(D)} & Modular characteristic
  & Curved Swiss-cheese over $\overline{\mathcal{M}}_g$
    (\S\ref{sec:bar_cobar}) \\[3pt]
\textbf{(H)} & Hochschild ring
  & Bulk $\simeq$ chiral Hochschild; BV-BRST origin
    (\S\ref{sec:chiral_hochschild_expanded})
\end{tabular}
\end{center}


\section*{What the bar complex computes}

Three descent mechanisms extract the invariants encoded in the
bar complex.

\emph{Cohomological descent} produces the $(-1)$-shifted Poisson
vertex algebra on $H^\bullet(\cA, Q)$
(Theorem~\ref{thm:cohomology_PVA}).  The higher operations
$m_{k \geq 3}$ vanish on cohomology via homotopies from the
contractibility of the topological factor in
$\FM_k(\C) \times \Conf_k(\R)$; the binary operation~$m_2$
decomposes into commutative product (regular part) and
$\lambda$-bracket (singular part); the PVA axioms---Jacobi,
Leibniz, skewsymmetry---follow from Arnold relations and Stokes'
theorem on $\FM_3(\C)$, the same identities that make
$d_{\barB}^2 = 0$ in Volume~I.

\emph{Spectral descent} produces the $R$-matrix from the
$\lambda$-bracket.  The bridge is the Laplace transform
\begin{equation}
\label{eq:intro-r-matrix}
r(z) \;=\; \int_0^\infty e^{-\lambda z}\,
\{{\cdot}\,{}_\lambda\,{\cdot}\}\, d\lambda,
\end{equation}
converting the PVA datum into a Yang--Baxter datum.  The quantum
$R$-matrix $R(z)$ from bulk-boundary composition
(Section~\ref{sec:spectral_braiding}) deformation-quantizes this
classical limit.

\emph{Genus descent} produces the curved Swiss-cheese algebras
over~$\overline{\mathcal{M}}_g$.  The curvature
$d^2 = \kappa(\cA) \cdot \omega_g$ couples the
$\R$-factorization to the Hodge bundle; the non-vanishing of
higher $\Ainf$ operations at the chain level IS this curved bar
structure---formality fails precisely because the categorical
logarithm acquires monodromy.

The three descents converge on the \emph{bulk-boundary-line
Koszul triangle}: the bulk chiral algebra
$\cA_{\mathrm{bulk}}$ on $\C \times \R$; the boundary vertex
algebra $\cA^!$ (Koszul dual) on~$\R$; the line category
$\mathcal{C}_{\mathrm{line}} \simeq
\cA^!\text{-}\mathbf{mod}$
(Theorem~\ref{thm:lines_as_modules}).  The three vertices are
related by
\[
A_{\mathrm{bulk}} \;\simeq\;
Z_{\mathrm{der}}(B_\partial) \;\simeq\;
\mathrm{HH}^\bullet(\cA^!)\text{:}
\]
the bulk is the \emph{derived center} of the boundary.  This is
the mathematical content of the holographic principle in the
holomorphic--topological setting.

Beyond the Heisenberg atom, three families test the full
machine.  The Landau--Ginzburg model provides the first
\emph{interacting} example: a genuine $m_3$ from a cubic
superpotential, truncating by degree.  Abelian Chern--Simons
provides the first \emph{braided} example, with an explicit
spectral $R$-matrix from the CS propagator and boundary
$\widehat{\mathfrak{u}(1)}_k$.  The Virasoro and $\mathcal{W}_3$
algebras connect to Volume~I's Drinfeld--Sokolov framework, with
recursive $m_k$ from the BV master equation and classical
Yang--Baxter from the $\lambda$-bracket Jacobi identity.


\section*{Reading this volume}

\noindent\textbf{Three entry points.}
\begin{enumerate}[label=\textup{(\alph*)}]
\item \emph{Via the Heisenberg atom.}
  Rosetta Stone (\S\ref{sec:rosetta-stone}) $\to$
  Foundations (\S\ref{sec:foundations}) $\to$
  PVA descent (\S\ref{sec:PVA_descent}) $\to$
  Concordance (\S\ref{sec:concordance}).
\item \emph{Complete theory.}
  Parts~I--II linearly, then examples (Part~IV), then
  Parts~III and~V.
\item \emph{Physics entry.}
  BV construction (\S\ref{sec:BV_construction}) $\to$
  FM calculus (\S\ref{sec:FM_calculus}) $\to$
  Examples (\S\ref{sec:examples_complete}) $\to$
  YM synthesis (\S\ref{sec:ym-synthesis}).
\end{enumerate}
Part~I builds the Swiss-cheese identification.  Part~II develops
the descent calculus.  Part~III develops the Koszul triangle and
its consequences.  Part~IV tests the machine on five families.
Part~V pushes the frontier: quantization, Yang--Mills, celestial
holography, anomaly completion.  The examples of Part~IV are
independent of Parts~III and~V.
